% =============================================================================
% Capítulo 6 — Conclusões
% =============================================================================
\section{Resumo dos achados}

A \textit{pipeline} desenvolvida neste trabalho demonstrou que a detecção automatizada de ocultações estelares em curvas de luz pode atingir desempenho elevado --- F1-score entre 0,978 e 0,994 e AUC-ROC superior a 0,998 --- utilizando \textit{features} clássicas extraídas das séries temporais e \textit{ensembles} de árvores de decisão (Random Forest, XGBoost, CatBoost), além de Regressão Logística como \textit{baseline}.

Os seis experimentos realizados permitiram identificar os seguintes padrões:

\begin{enumerate}
	\item \ia{\textbf{O \textit{feature engineering} baseado em domínio é o fator determinante da qualidade da classificação.}} A redução de 28 para 13--14 \textit{features} (remoção de variáveis redundantes) não degradou --- e em alguns casos melhorou --- o desempenho, indicando que um conjunto compacto de variáveis fisicamente motivadas é suficiente.

	\item \ia{\textbf{A \textit{feature} \texttt{kmeans\_centroid\_dist} é dispensável apesar de sua alta importância atribuída pelos modelos de árvore.}} O estudo de ablação (Experimento~6) mostrou que a remoção dessa variável causa degradação desprezível ($\leq$0,3~pp em F1), demonstrando que importância de \textit{feature} não é sinônimo de insubstituibilidade: outras variáveis colineares (sobretudo \texttt{Feature\_Savgol\_std}) compensam sua ausência.

	\item \ia{\textbf{O ajuste do limiar de decisão permite operar a \textit{pipeline} em regime de alta sensibilidade}} (minimização de falsos negativos) sem retreinamento, explorando a separabilidade elevada entre as classes. O estudo de caso de Quaoar (Seção~\ref{sec:estudo_caso_quaoar}) demonstrou, em dados reais, que essa alavanca recupera anéis finos que passariam despercebidos no limiar padrão.
\end{enumerate}

Como verificação complementar, as 67 curvas do banco com $\mathrm{Occ\_SNR\_dip} < 3$ foram submetidas aos modelos treinados; todas foram corretamente classificadas como positivas pelos três \textit{ensembles} de árvores (67/67), e 66 das 67 pela Regressão Logística. Embora essas curvas tenham participado do treinamento e o teste não inclua negativos de baixo SNR, o resultado indica que os modelos não descartam sistematicamente eventos ruidosos --- condição necessária (embora não suficiente) para a aplicabilidade da \textit{pipeline} em cenários de baixa relação sinal-ruído.

% -----------------------------------------------------------------------------
\section{Resposta à hipótese de pesquisa}

A hipótese formulada na Introdução --- de que é possível obter desempenho comparável a abordagens baseadas em redes convolucionais ou inspeção manual usando \textit{features} clássicas e \textit{ensembles} de árvores, com vantagem de interpretabilidade --- é \textbf{corroborada} pelos resultados obtidos. Os quatro modelos testados atingiram desempenho \ia{elevado e consistente}; a concordância entre modelos de naturezas distintas (linear vs.\ árvores, \textit{bagging} vs.\ \textit{boosting}) quanto à importância das variáveis reforça a interpretabilidade do sistema e a solidez das conclusões.

% -----------------------------------------------------------------------------
\section{Contribuições originais}

As principais contribuições deste trabalho são:

\begin{enumerate}
	\item \textbf{Pipeline completa de detecção:} desde a ingestão de curvas de luz já reduzidas até a classificação binária com reportagem de métricas, importância de variáveis e predições por curva, implementada em Python e disponível para reprodução.

	\item \textbf{Conjunto de \textit{features} baseado em domínio:} 13--14 variáveis que capturam assinaturas físicas de ocultações (estatísticas de Savitzky-Golay, \textit{Max Drawdown}, métricas observacionais de profundidade e SNR, testes estatísticos entre quartis), selecionadas por análise de redundância e validadas experimentalmente.

	\item \textbf{Estudo de ablação da \textit{feature} K-Means:} \ia{demonstração empírica de que a importância atribuída por algoritmos de árvore não implica insubstituibilidade, com implicação prática para simplificação do \textit{pipeline}.} \ia{[Repete o achado n.º 2 acima, na mesma página --- unificar os dois itens.]}

	\item \textbf{Análise de sensibilidade ao volume de treino:} evidência de saturação da curva de aprendizado, orientando futuras decisões sobre coleta de dados.

	\item \textbf{Análise de ajuste de limiar:} demonstração de que a assimetria de custo entre falsos negativos e falsos positivos pode ser incorporada em pós-processamento, sem retreinamento.

	\item \textbf{Validação cruzada estratificada e teste de McNemar:} estimativa robusta da variabilidade do desempenho (desvio padrão pequeno entre dobras) e avaliação da significância estatística das diferenças entre modelos, que se mostraram equivalentes no conjunto de teste.

	\item \textbf{Demonstração da capacidade de (re)visitar curvas em busca de eventos sutis:} o estudo de caso de Quaoar (Seção~\ref{sec:estudo_caso_quaoar}) mostrou, em dados reais externos ao treino, que o modelo separa um anel fino verdadeiro de um trecho de ruído semelhante por quase duas ordens de grandeza de probabilidade --- contraste que, via ajuste de limiar, permite recuperar eventos que escapariam à classificação binária padrão.
\end{enumerate}

% -----------------------------------------------------------------------------
\section{Limitações}

As limitações deste trabalho devem ser explicitadas para contextualizar adequadamente os resultados:

\begin{enumerate}
	\item \textbf{\textit{Dataset} parcialmente sintético:} das 891 curvas negativas, 702 são sintéticas e 186 são negativos por recorte (segmentos reais extraídos de curvas positivas); apenas 3 são detecções negativas nativas. Embora as curvas sintéticas tenham sido geradas com base em propriedades estatísticas de curvas reais, a capacidade de generalização para um cenário puramente observacional não foi validada.

	\item \textbf{\textit{Split} único por experimento:} cada experimento utiliza uma única partição treino/teste. A implementação de validação cruzada estratificada (Seção~5.\ref{sec:cross_validation}) mitiga parcialmente essa limitação; uma validação cruzada repetida com múltiplas sementes fortaleceria as conclusões.

	\item \textbf{Ausência de validação em catálogos externos:} as curvas reais utilizadas neste trabalho provêm de duas fontes distintas: o catálogo \texttt{B/occ/asteroid} do VizieR - agregador público que reúne campanhas independentes envolvendo múltiplos observadores, telescópios e corpos - e o banco do Grupo do Rio (campanhas específicas de Umbriel e Chiron). Embora a heterogeneidade interna do VizieR já garanta diversidade observacional considerável, a generalização da \textit{pipeline} para curvas oriundas de outros agregadores ou de levantamentos específicos não foi avaliada sistematicamente. Apesar disso, o script utilizado para popular o banco SQL com dados do Grupo do Rio serve para qualquer tipo de curvas de luz, contanto que estejam em formato \texttt{.dat}.

	\item \textbf{Hiperparâmetros fixos:} os hiperparâmetros dos modelos foram definidos manualmente com base em valores recomendados pela literatura; não foi realizada busca em grade ou otimização bayesiana.

	\item \textbf{Predominância de curvas com ocultações bem definidas:} o elevado desempenho dos modelos reflete, em parte, a predominância de curvas com ocultações claras (alto S/N, queda pronunciada) no conjunto de dados. Para curvas com baixo S/N ou eventos rasantes, a tarefa de classificação é substancialmente mais difícil e o desempenho real da \textit{pipeline} nessas condições ainda não foi avaliado sistematicamente.
\end{enumerate}

% -----------------------------------------------------------------------------
\section{Trabalhos futuros}

Com base nos resultados e limitações identificados, sugerem-se as seguintes direções para trabalhos futuros:

\begin{itemize}
	\item \textbf{Validação cruzada repetida:} executar a validação cruzada com múltiplas sementes (e.g., 10 repetições $\times$ 5 dobras) para obter intervalos de confiança robustos das métricas.

	\item \textbf{Experimento \textit{real holdout}:} reservar todas as curvas reais (positivas e negativas genuínas) para o conjunto de teste, treinando apenas com curvas sintéticas e negativos por recorte, para avaliar a transferibilidade do modelo.

	\item \textbf{\textit{Features} baseadas em difração de Fresnel:} incorporar características derivadas de modelos físicos mais sofisticados (escala de Fresnel, opacidade parcial, diâmetro aparente da estrela), que poderiam melhorar a discriminação em eventos de baixo SNR.

	\item \textbf{Detecção em regime de baixo S/N:} construir um subconjunto de teste composto exclusivamente por curvas difíceis --- baixo S/N, ocultações rasantes, eventos com poucos pontos no \textit{dip} --- para avaliar o desempenho da \textit{pipeline} nas condições em que a triagem automática é de fato necessária, isto é, onde a inspeção visual é ambígua.

	\item \textbf{Janela deslizante e \textit{features} locais:} o estudo de recortes de Quaoar (Seção~\ref{sec:estudo_caso_quaoar}) mostrou que o ajuste de limiar resolve o caso prático de detecção de anéis finos, mas que a causa estrutural da baixa probabilidade absoluta atribuída a esses eventos é o caráter \textbf{global} das \textit{features} (estatísticas agregadas sobre toda a janela diluem um \textit{dip} de poucos segundos). Duas direções complementares atacam essa raiz: (i)~aplicar a \textit{pipeline} por \textbf{janela deslizante} de largura compatível com a duração típica do evento alvo, tornando o conjunto atual de \textit{features} sensível a eventos locais; e (ii)~enriquecer o conjunto com \textbf{\textit{features} locais} --- profundidade máxima de qualquer subjanela de $N$ pontos abaixo do \textit{baseline}, simetria do perfil de ingresso/egresso, ajuste a poço retangular curto.

	\item \textbf{Classificação multi-classe de eventos:} estender a formulação binária (ocultação presente/ausente) para uma classificação multi-classe que distinga o \textbf{tipo} de evento observado: ocultação por corpo sólido, assinatura atmosférica (refração e absorção), detecção de anéis e padrões de difração de Fresnel. Essa extensão exigiria \textit{features} adicionais --- por exemplo, simetria do perfil de ingresso/egresso, profundidade relativa de quedas secundárias e ajuste a modelos de difração --- além de um \textit{dataset} curado com rótulos para cada categoria, possivelmente com apoio de curvas sintéticas geradas por simuladores que incluam esses fenômenos.
\end{itemize}
